% =============================================================================
% commands.tex —— 目录、脚注等「分散在各处的小设置」
% -----------------------------------------------------------------------------
% 依赖：base.tex（\ink@sourcenotesize 等）、colors.tex、layout.tex。
% 页码相关的 \ink@pagenumber 定义在 layout.tex（与页脚模板放在一起）。
% =============================================================================

% -----------------------------------------------------------------------------
% 1. 目录（两级）字体
% -----------------------------------------------------------------------------
% 一级目录：正文尺寸 + 衬线粗体；二级目录：小字号 + 无衬线常规字重。
% 具体字号在style.tex 中覆盖（小版 \scriptsize，大版 \footnotesize），
% 这里只固定字族与字形，保证两版一致。
\setbeamerfont{section in toc}{size=\normalsize,family=\rmfamily,series=\bfseries,shape=\upshape}
\setbeamerfont{subsection in toc}{family=\sffamily,series=\mdseries,shape=\upshape}

% -----------------------------------------------------------------------------
% 2. 脚注
% -----------------------------------------------------------------------------
% 脚注在学术汇报里通常承载引文（作者、标题、会议、年份），因此正文与粗体
% 都使用 Times 风格字体（\ink@citefont，定义在 fonts.tex），字形保持 Italic；
% 脚注编号保持直立（斜体数字难辨认）。
% 字号由style.tex 设置（小版 \scriptsize，大版 \footnotesize）。
\setbeamerfont{footnote}{family=\ink@citefont,shape=\itshape}
\setbeamerfont{footnote mark}{shape=\upshape}
